\renewcommand{\ElicitationDate}{2026.10.29}

\section{Elicitation \ElicitationDate}

\begin{tabular}{ll}
\textbf{Date:} & \ElicitationDate \\
\textbf{Consultant:} & Hector Hill / Gy'ax \\
\textbf{Elicitors:} & Zahra, Andrew, Matthias \\
\textbf{Location:} & TFS elicitation room \\
\textbf{Context:} & Sixth elicitation for LING 431/531 Field Methods \\
\textbf{Audio:} & \texttt{October\_29\_2026\_HH.WAV} \\
\end{tabular}

\subsection{Matthias [0:00-34:00]}

[1:00]: \textit{ligi wil} is a gift for feasts and such... not sure if this is composed of \Dwid{} and \Comp{}. Could be related to \textit{li'ligit} `feast'?

\textit{sax̱ gi'namaa} is a gift that you give without expectations to receive something back. 

[2:10]:

\begin{exe}
    \ex \glll gi'nam'y sax̱ gi'namaa loon \\
    gi'nam-'y sax̱ gi'nam-aa loo-n  \\
    give-\Fsg{}.\Ii{} only give-\Nmlz{}? \Obl{}-\Ssg{}.\Ii{} \\
    `I gave a gift to you.'
\end{exe}

[3:18]: I should confirm if \textit{gat} here can mean ``person'' or only ``man'' specifically.

\begin{exe}
    \ex \glll gi'nam'y sax̱ gi'namaa ahl bag̱adil(hl) gat \\
    gi'nam-'y sax̱ gi'nam-aa a[-t]=hl bag̱adil(=hl) gat \\
    give-\Fsg{}.\Ii{} only give-\Nmlz{}? \Prep{}[-3.\Ii{}]=\Cn{} two.\Hum{}(=\Cn{}) person \\
    `I gave a gift to two people.'
\end{exe}

[3:58]: at [5:45] HH says ``\textit{hilda} doesn't add on to the \textit{gat}''. It's kind of curious that he doesn't accept \textit{hilda} here, since Zahra had a few examples with \textit{hilda hun} where it was accepted.

\begin{exe}
    \ex 
    \begin{xlist}
    \item \glll * gi'nam'y sax̱ gi'namaa ahl hilda gat \\
    {} gi'nam-'y sax̱ gi'nam-aa a[-t]=hl hilt-a gat \\
    {} give-\Fsg{}.\Ii{} only give-\Nmlz{}? \Prep{}[-3.\Ii{}]=\Cn{} many-\Attr{} person \\
    Intended: `I gave a gift to many people.'
    \item \glll * gi'nam'y sax̱ gi'namaa ahl hilthl gat \\
    {} gi'nam-'y sax̱ gi'nam-aa a[-t]=hl hilt=hl gat \\
    {} give-\Fsg{}.\Ii{} only give-\Nmlz{}? \Prep{}[-\Ssg{}.\Ii{}]=\Cn{} many=\Cn{} person \\
    Intended: `I gave a gift to many people.'
    \end{xlist}
\end{exe}

[6:32]: Andrew provided this form, and Hector techincally didn't reject it, but its status is inconclusive. Only comment is that ``that's a question''.

\begin{exe}
    \ex \glll gi'nam'y sax̱ gi'namaa ahl hilt gataa \\
    gi'nam-'y sax̱ gi'nam-aa a[-t]=hl hilt gat=aa \\
    give-\Fsg{}.\Ii{} only give-\Nmlz{}? \Prep{}[-\Ssg{}.\Ii{}]=\Cn{} many person=\Q{} \\
    \textit{Grammaticality inconclusive: }`Did you give a gift to many people?'
\end{exe}

[7:08]: But the minimal pair of [6:32] without the \Q{} is rejected.

\begin{exe}
    \ex \glll * gi'nam'y sax̱ gi'namaa ahl hilt gat \\
    {} gi'nam-'y sax̱ gi'nam-aa a[-t]=hl hilt gat \\
    {} give-\Fsg{}.\Ii{} only give-\Nmlz{}? \Prep{}[-\Ssg{}.\Ii{}]=\Cn{} many person \\
    Intended: `I gave a gift to many people.'
\end{exe}

[7:27]: Here HH is attempting to translate the sentence with \textit{hilthl} but instead uses \textit{walḵ'a 'nithl}. Obviously this is not the same meaning.

It is also weird that the direct object has a \Prep{} before it, but that is how it sounds. Unless this is a production error from the more deliberate translation task.

This sentence without \textit{ahl} is repeated at [9:53] and [11:02]. But, at [8:05] there is a translation with the direct object being preceded by \textit{ahl} as well. Is it optional here? This is my tentative gloss.

\begin{exe}
    \ex \glll sax̱ gi'nam'y (ahl) sax̱ gi'namaa ahl walḵ'a 'nithl gat \\
    sax̱ gi'nam-'y (a[-t]=hl) sax̱ gi'nam-aa a[-t]=hl walḵ'a 'nit=hl gat \\
    only give-\Fsg{}.\Ii{} (\Prep{}[-\Ssg{}.\Ii{}]=\Cn{}) only give-\Nmlz{}? \Prep{}[-\Ssg{}.\Ii{}]=\Cn{} all 3.\Iii{}=\Cn{} person \\
    `I gave a gift to all the people.' 
\end{exe}

[8:05]:

\begin{exe}
    \ex \glll sax̱ gi'nam'y ahl hildhl sax̱ gi'namaa ahl walḵ'a 'nithl gat \\
    sax̱ gi'nam-'y a[-t]=hl hild=hl sax̱ gi'nam-aa a[-t]=hl walḵ'a 'nit=hl gat \\
    only give-\Fsg{}.\Ii{} \Prep{}[-\Ssg{}.\Ii{}]=\Cn{} many=\Cn{} only give-\Nmlz{}? \Prep{}[-\Ssg{}.\Ii{}]=\Cn{} all 3.\Iii{}=\Cn{} person \\
    \textit{Inconclusive: }`I gave a lot of gifts to all the people.' 
\end{exe}

[11:30]: This was also provided when asking to translate ``I gave a gift to \textit{many} people''.

\begin{exe}
    \ex \glll gi'nam'y  sax̱ gi'namaa ahl gat \\
    gi'nam-'y  sax̱ gi'nam-aa a[-t]=hl gat \\
    give-\Fsg{}.\Ii{} only give-\Nmlz{}? \Prep{}[-\Ssg{}.\Ii{}]=\Cn{} person \\
    `I gave a gift to (some) people.' 
\end{exe}

[12:05]: ``If you want to limit the \textit{gat}, you can say'': 

This still doesn't get at the meaning of ``many men'' or ``many people'' unfortunately, but it's a step closer.


\begin{exe}
    \ex \glll gi'nam'y  sax̱ gi'namaa ahl walḵ'a 'nithl gat, ii nee diin gi'namt ahl hanaḵ' \\
    gi'nam-'y  sax̱ gi'nam-aa a[-t]=hl walḵ'a 'nit=hl gat, ii nee dii=n gi'nam-t a[-t]=hl hanaḵ' \\
    give-\Fsg{}.\Ii{} only give-\Nmlz{}? \Prep{}[-\Ssg{}.\Ii{}]=\Cn{} all 3.\Iii{}=\Cn{} person, \Ccnj{} \Neg{} \Foc{}=\Fsg{}.\I{} give-\Tsg{}.\Ii{} \Prep{}[-\Ssg{}.\Ii{}]=\Cn{} woman \\
    `I gave a gift to all the men and none of the women.' 
\end{exe}

[14:01]: HH translated \textit{hlag̱ats'uuhl gat} as ``some of the people''.

I could not find \textit{sida/seda} in any resource other than FV (\url{https://www.firstvoices.com/gigeenix-gitxsanimx/words/0df4a112-cad3-4d4a-8f9b-d52c68874565}). I'm assuming it's a preverb.

Also note that HH clearly says that this is repitive, in that you would say this for ``people who don't understand''.

\begin{exe}
    \ex \glll gi'nam'y sax̱ gi'namaa ahl hlag̱ats'uuhl gat, nee diin sida gi'namhl sax̱ gi'namaa ahl walḵ'a 'nithl gat \\
    gi'nam-'y  sax̱ gi'nam-aa a[-t]=hl hla-g̱a-ts'uu=hl gat, nee dii=n sida gi'nam=hl sax̱ gi'nam-aa a[-t]=hl walḵ'a 'nit=hl gat  \\
    give-\Fsg{}.\Ii{} only give-\Nmlz{}? \Prep{}[-\Ssg{}.\Ii{}]=\Cn{} \Part{}-\Distr{}-other=\Cn{} person, \Neg{} \Foc{}=\Fsg{}.\I{} all give-\Cn{} only give-\Nmlz{}? \Prep{}[-\Ssg{}.\Ii{}]=\Cn{} all 3.\I{}=\Cn{} person \\
    `I gave a gift to some people, and I didn't give a gift to all the people.' 
\end{exe}

[15:20]:

\begin{exe}
    \ex \glll nee diin gi'namhl sax̱ gi'namaa ahl walḵ'a 'nithl gat, iin k'ap gi'namt ahl  hlag̱ats'uut \\
    nee dii=n gi'nam=hl sax̱ gi'nam-aa a[-t]=hl walḵ'a 'nit=hl gat, ii=n k'ap gi'nam-t a[-t]=hl hla-g̱a-ts'uu-t \\
    \Neg{} \Foc{}=\Fsg{}.\I{} give=\Cn{} only give-\Nmlz{}? \Obl{}[-3.\Ii{}]=\Cn{} all 3.\I{}=\Cn{} person, \Ccnj{}=\Fsg{}.\I{} \Verum{} give-\Tsg{}.\Ii{} \Obl{}[-3.\Ii{}]=\Cn{} \Part{}-\Distr{}-other-\Tsg{}.\Ii{} \\
    `I didn't give a gift to all the people, I gave it to the other ones.' 
\end{exe}

[16:33]: 

\begin{exe}
    \ex 
    \begin{xlist}
        \item \glll anoogi'y  sax̱ gi'namaahl bag̱adilhl gat \\
    anoogi-'y  sax̱ gi'nam-aa=hl bag̱adil=hl gat  \\
    like-\Fsg{}.\Ii{} only give-\Nmlz{}?=\Cn{} two.\Hum{}=\Cn{} person \\
    `I liked two people's gifts.' 
        \item \glll sim lukw'il anoogi'yhl  sax̱ gi'namaahl bag̱adilhl gat \\
    sim lukw'il anoogi-'y=hl  sax̱ gi'nam-aa=hl bag̱adil=hl gat  \\
    true very like-\Fsg{}.\Ii{}=\Cn{} only give-\Nmlz{}?=\Cn{} two.\Hum{}=\Cn{} person \\
    `I really liked two people's gifts.' 
    \end{xlist}
\end{exe}

[17:20]: Again, \textit{hilda} rejected here. ``that \textit{hilda}, `'\textit{hildaya}' that's a question, `are they a lot', `is there a lot'''

\begin{exe}
    \ex \glll * sim lukw'il anoogi'y  sax̱ gi'namaahl hilda gat \\
    {} sim lukw'il anoogi-'y  sax̱ gi'nam-aa=hl hilt-a gat  \\
    {} true very like-\Fsg{}.\Ii{} only give-\Nmlz{}?=\Cn{} many-\Attr{} person \\
    Intended: `I really liked many people's gifts.' 
\end{exe}

[18:35]: I'm pretty sure the \textit{=hl} here is combining to the noun, but it sounds a bit like \textit{a[-t]=hl}.

\begin{exe}
    \ex \glll sim lukw'il anoogi'yhl  sax̱ gi'namaahl (walḵ'a 'nithl) gat \\
    sim lukw'il anoogi-'y=hl  sax̱ gi'nam-aa=hl (walḵ'a 'nit=hl)  gat  \\
    true very like-\Fsg{}.\Ii{} only give-\Nmlz{}?=\Cn{} (all 3.\I{}=\Cn{}) person \\
    `I really liked (all the) people's gifts.' 
\end{exe}

[21:10]:

\begin{exe}
    \ex \glll hilthl gat ant sax̱ gi'namaahl anoogi'y \\
    hilt=hl gat an=t sax̱ gi'nam-aa=hl anoogi-'y  \\
    many=\Cn{} people \Ax{}=3.\I{} only give-\Nmlz{}?=\Cn{} like-\Fsg{}.\Ii{} \\
    `I liked a lot of people's gifts.'  Lit: `The gifts of people that I like are many.'
\end{exe}

[21:50]:

\begin{exe}
    \ex \glll sim anoogi'y walḵ'a 'nithl sax̱ gi'namaahl gat \\
    hilt=hl gat an=t sax̱ gi'nam-aa=hl anoogi-'y  \\
    many=\Cn{} people \Ax{}=3.\I{} only give-\Nmlz{}?=\Cn{} like-\Fsg{}.\Ii{} \\
    `I like many people's gifts.'
\end{exe}

[22:08]: \textit{hilt} can also fit in like this it looks like. 

\begin{exe}
    \ex \glll sim anoogi'y hilthl sax̱ gi'namaahl gat \\
    sim anoogi-'y hilt=hl sax̱ gi'nam-aa=hl gat \\
    true like-\Fsg{}.\Ii{} many=\Cn{} only give-\Nmlz{}?=\Cn{} person \\
    `I like people's many gifts' Lit: `I like people's gifts that are many.'
\end{exe}

[22:56]: We can also break it up like this. HH explicitly mentions putting the commas where they are.

\begin{exe}
    \ex \glll hilthl gat, sax̱ gi'namaadiit, anoogi'y \\
    hilt=hl gat, sax̱ gi'nam-aa-diit, anoogi-'y\\
    many=\Cn{} person, only give-\Nmlz{}?-\Tpl{}.\Ii{}, like-\Fsg{}.\Ii{} \\
    `Many people, their gifts, I like.'
\end{exe}

[24:25]: I'm not actually sure if something like a comma is required in b.

\begin{exe}
    \ex 
    \begin{xlist}
        \item \glll ga'a'y gilbil smex (yukwhl) x̱maa'ydiit \\
    ga'a-'y gilbil smex (yukw=hl) x̱-maa'y-diit \\
    see-\Fsg{}.\Ii{} two bear (\Prog{}=\Cn{}) \Expr{}-berry-\Tpl{}.\Ii{} \\
    `I saw two bears eating berries'
    \item \glll x̱maa'yhl gilbil smex ga'a'y \\
    x̱-maa'y=hl gilbil smex ga'a-'y \\
    \Expr{}-berry=\Cn{} two bear see-\Fsg{}.\Ii{} \\
    `Two bears eating bearies, I saw'
    \end{xlist}
\end{exe}

[25:25]: Sentence confirmed by Hector. I didn't confirm the use of the comma though, so that is tentative.

\begin{exe}
    \ex \glll gilbil smex, gya'a'y yukwhl \tsb{xmaa'ydiit} \\
    gilbil smex, gya'a-'y yukw=hl x̱-maa'y-diit \\
    two bear, see-\Fsg{}.\Ii{} \Prog{}=\Cn{} \Expr{}-berry-\Tpl{}.\Ii{} \\
    `Two bears, I saw them eating berries.'
\end{exe}

[25:40]:

\begin{exe}
    \ex \glll hilthl smex, gya'a'y yukwhl \tsb{xmaa'ydiit} \\
    hil=hl smex, gya'a-'y yukw=hl x̱-maa'y-diit \\
    many=\Cn{} bear, see-\Fsg{}.\Ii{} \Prog{}=\Cn{} \Expr{}-berry-\Tpl{}.\Ii{} \\
    `Many bears, I saw them eating berries.'
\end{exe}

[25:55]: ``Broken but it's understandable.'

\begin{exe}
    \ex \glll * gya'a'y hilthl smex yukwhl \tsb{xmaa'ydiit} \\
    {} gya'a-'y hil=hl smex  yukw=hlx̱-maa'y-diit \\
    {} see-\Fsg{}.\Ii{} many=\Cn{} bear \Prog{}=\Cn{} \Expr{}-berry-\Tpl{}.\Ii{} \\
    Intended: `I saw many bears eating berries.'
\end{exe}


[27:03]:

\begin{exe}
    \ex \glll gya'a'y win guphl smexhl maa'y \\
    gya'a-'y win gup=hl smex=hl maa'y \\
    see-\Fsg{}.\Ii{} \Comp{} eat=\Cn{} bear=\Cn{} berry  \\
    `I saw bears eat berries.'
\end{exe}

[28:18]: Comma specified by HH. Translation is my best approximation.

\begin{exe}
    \ex \glll gya'a'y win tooḵwhl smex, yukwhl gupdiit maa'y  \\
    gya'a-'y win tooḵw=hl smex, yukw=hl gup-diit maa'y  \\
    see-\Fsg{}.\Ii{} \Comp{} eat.\Pl{}=\Cn{} bear, \Prog{}=\Cn{} eat-\Tpl{}.\Ii{} berry  \\
    `I saw the bears eating, they were eating berries'
\end{exe}

[29:05]: \textbf{Context: }Telling people how many people will be at a party. (Re-elicitation of 10.08)

Negative judgement here is confirmed at [30:15].


\begin{exe}
    \ex \glll dim hilt(*hl) 'nuu'm  \\
    dim hilt(*=hl) 'nuu'm \\
    \Prosp{} many(*=\Cn{}) \Fpl{}.\Iii{}  \\
    `We will be many.'
\end{exe}

[29:50]: Reminder here that we still don't know what \textit{dimis/da'mis} is. Above is a good answer to this.

\begin{exe}
    \ex
    \begin{xlist}
    \item \glll dimis dim g̱abisi'm \\
    dimis dim g̱abi-si'm \\
    dimis \Prosp{} \Cnt{}-\Spl{}.\Ii{} \\
    `How many will you be?'
    \glll nde dim g̱abisi'm \\
    nde dim g̱abi-si'm \\
    \Wh{} \Prosp{} \Cnt{}-\Spl{}.\Ii{} \\
    `How many will you be?'
    \end{xlist}
\end{exe}


[31:50]:

\begin{exe}
    \ex \glll dim hilt 'nisi'maa  \\
    dim hilt 'nisi'm=aa \\
    \Prosp{} many \Spl{}.\Iii{}=\Q{}  \\ 
    `Are there gonna be a lot of you guys?' \textit{VF}
\end{exe}

[32:20]: Also in a question-response context: 

Also assuming that b. doesn't actually get a prospective reading, but moreso a present tense reading. At [32:25] HH translates it to present in English.

Also, at [33:00] ``\textit{xwsidns 'nuu'm} is like there are five of these that you are gonna bring each.'' Also, HH says that this might fit in Gigeenix since he's heard that before, and may also fit with Gitanyow and Kispiox.

\begin{exe}
    \ex
    \begin{xlist}
        \item \glll * dim xwsdins 'nuu'm  \\
    {} dim xwsdins 'nuu'm \\
    {} \Prosp{} five \Fpl{}.\Iii{} \\
    Intended: `We will be many.'
    \item xwsdinshl g̱abi'm  \\
    xwsdins=hl g̱abi-'m \\
    five=\Cn{} \Cnt{}-\Fpl{}.\Ii{}  \\ 
    `There are five of us.'
    \end{xlist}
\end{exe}

\subsection{Andrew [34:00-57:45]}


[34:41] --> when re-testing the claim that sag̱ayt cannot be self-reflexive from 10/15 [16:25], Hector seems to change his mind. To be fair there is much less context here
\begin{exe}
\ex \gll k'i=hl g̱an=hl k'i ii sag̱ayt gilbil-t\\
one=\Cn{} \Pcnj{}=\Cn{} one and all.together two-3.\Ii{}\\
'One and one is two'
\end{exe}
HH: \textit{"gilbilt is two"}

[35:10] --> a. is re-eliciting 10/15 [23:30] (sf); b. is Hector's vf
\begin{exe}
\ex \begin{xlist}
\item \gll k'i g̱an=hl k'i an-jep=hl gilbil\\
one \Pcnj{}=\Cn{} one ??-make=\Cn{} two\\
'one and one make two'
\item \gll k'i \tsb{gan}-t k'i e-jep=hl gilbil\\
one \Pcnj{}-3.\I{} one ??-make=\Cn{} two\\
\end{xlist}
\end{exe}

different w/r/t the g̱an suffix and the pre-verb, but that could be me not hearing 10/15's properly

[36:00] --> AS: sf subbing in \textit{jehl} from the addition constructions.
\begin{exe}
    \ex \gll * gwile' jehl k'i e-jep=hl gilbil\\
   {} three lose.at one ??-make=\Cn{} two\\
    \textit{Intended: "three losing one makes two"}
\end{exe}

HH: \textit{"good, but all broken up"}. I'll take what I can get; I consider this an absolute win...

[38:30] --> vf
\begin{exe}
\ex \gll k'i saa-guu-t-'y e=hl gwile'l-t ii ḵ'am gilbil=hl gina ment\\
one away-take-\Trn{}-1.\Ii{} \Prep{}=\Cn{} three-3.\Ii{} and only two=\Cn{} behind leftover\\
'I took one out three and now I have two'
\end{exe}

Very happy with this construction as it reinforces that \textit{ḵ'am} has to be in the difference phrase and it decomposes very well!

[40:45] --> sf
\begin{exe}
\ex \begin{xlist}
\item \gll g̱al hilt 'nuu'm\\
too many 3Pl.\Iii{}\\
\textit{Intended: "There's too many of us"}
\item \gll Gabii=hl g̱al hilt?\\
\Mass{}=\Cn{} too many\\
'How many too many?'
\item \gll k'i g̱al hilt 'nuu'm\\
one too many 3Pl.\Iii{}\\
\textit{Intended: "There's one too many of us"}
\end{xlist}
\end{exe}

[42:00] --> vf
\begin{exe}
\ex \gll k'i, win g̱al hilt-'m\\
one, \Comp{} too many-1Pl.\Iii{}\\
'one we're too many'
\end{exe}

HH: "There has to be a comma there, in front of \textit{k'i}"

[43:05]
\begin{exe}
\ex \gll * g̱al hilt-'m win k'i\\
{} too many-1Pl.\Iii{} \Comp{} one\\
\textit{Intended: we're one too many}
\end{exe}

HH: "broken"

[43:52]
\begin{exe}
\ex \gll hlibuu\\
few\\
\end{exe}

[45:29] --> sf
\begin{exe}
\ex \gll gilbil, win g̱al hlibuu=hl ḵawts\\
two, \Comp{} too few=\Cn{} carrots\\
'Two too few carrots'
\end{exe}

HH: "I like that you put, when you have numbers, you have to separate it; unless you have people, then it's a little bit different"

Investigate this claim further--how is it different?

[46:07] --> vf // HH: "You could say"
\begin{exe}
\ex \gll sk'okxw 'nuu e=hl ḵawts\\
not.enough 3Pl.\Iii{}(?) \Prep{}=\Cn{} carrots\\
'we don't have enough carrots'
\end{exe}

AS: Not sure about the orthography of \textit{sk'okxw}; that's the way I wrote it on the board during the elicitation, but FV and the storybook spell it differently, (\textit{sḵooksxw} and what looks similar \textit{sgoo=hl} [== so.many=CN] (p.212 of the .pdf) respectively). I bring this up because I'm not sure what the decomposition of the morphemes could be here. Might be worth asking.
This could also fully be a reach for the storybook comparison, since they're antonyms, being \textit{not enough} and \textit{so many}\dots

[47:05] --> sf
\begin{exe}
\ex \gll sgidim-'y gilbil hilt\\
should-1.\Ii{} two many\\
\textit{Intended: "I need two more"}
\end{exe}

As Matthias pointed out \textit{sgidim} is more modally and translates as \textit{should}, but it was the best I could bring into the elicitation. Thank you for the great context below:

[49:45] --> \textit{Context: You're going to the store and you ask "What do you need from the store"}
\begin{exe}
\ex \gll gu=hl asag̱a-n gaa e=hl g̱aldim woo't?\\
what=\Cn{} want-2.\Ii{} ?? \Prep{}=\Cn{} container sell \\
'what do you need at the store?'
\end{exe}

\textit{store} is a selling container; this is a very fun mental image, personally\dots

[51:25] --> sf
\begin{exe}
    \ex \gll * asag̱a-'y gilbil hilt\\
    {} want-1.\Ii{} two many\\
\textit{Intended: "I need two more"}
\end{exe}

[52:05] --> vf
\begin{exe}
\ex \gll he'ksimux asag̱a-'m gilbil=hl ḵawts\\
again want-1Pl.\Ii{} two=\Cn{} carrots\\
'We two more carrots'
\end{exe}

[52:49] --> sf
\begin{exe}
\ex \begin{xlist}
\item \gll \# he'ksimux asag̱a-'m legim anaax\\
{} again want-1Pl.\Ii{} flour bread\\
\textit{Intended: "I need more flour"}
\item \gll he'ksimux asag̱a-'m g̱aldim legim anaax\\
again want-1Pl.\Ii{} container flour bread\\
The difference being that stores sell bags of flour and not just a pile\dots makes sense
\item \gll he'ksimux asag̱a-'m dihlxw legim anaax\\
again want-1.Pl\Ii{} bag flour bread\\
'We need a bag of flour'
\end{xlist}
\end{exe}

[53:55]
\begin{exe}
\ex \gll g̱ay k'aa asag̱a-'m gilbil ḵawts \\
\Cntr{} exceedingly want-1.Pl.\Ii{} two carrots \\
\end{exe}

\subsection{Zahra [58:05-1:28:37]}

[58:]

\begin{exe}
    \ex \glll Gwile'lhl sise’ehl g̱a’nagwihl ha’niitx̱ooḵxw ii gilbilhl sise’ehl g̱ale’yhlit, ii sag̱ayt ḵ’oolt g̱asg̱oot. \\
    gwile'l=hl si-se’e=hl g̱a-’nakw=hl ha’niitx̱ooḵxw ii gilbil=hl si-se’e=hl   g̱a-le’y=hl-it, ii sag̱ayt ḵ’oolt g̱asg̱oo-t. \\
    three=\Cn{} \Pl-foot=\Cn{} \Distr-long=\Cn{} table and two \Pl-foot=\Cn{} \Distr-long=\Cn-3.\Ii, \Ccnj{} together six \Mass-3.\Ii \\ 
    \glt `The table is three feet long and two feet wide; together six [feet in area].'\hfill HH.\ElicitationDate{} (SF)
\end{exe}

We attempted to re-elicit the above (which had seemed more or less acceptable last time) but this time HH preferred \textit{xwsdins} rather than \textit{ḵ'oolt}, i.e. adding rather than multiplying. Andrew supplied a context with a 3x2 table setting with plates. HH: "If we're talking about the plates, this [the above] is right."

[1:03:23]
\begin{exe}
    \ex \glll Gwile'lhl sise’ehl g̱a’nagwihl ha’niitx̱ooḵxw ii gilbilhl sise’ehl g̱ale’yhlit, ii sag̱ayt xwsdins g̱asg̱oot. \\
    gwile'l=hl si-se’e=hl g̱a-’nakw=hl ha’niitx̱ooḵxw ii gilbil=hl si-se’e=hl   g̱a-le’y=hl-it, ii sag̱ayt xwsdins g̱asg̱oo-t. \\
    three=\Cn{} \Pl-foot=\Cn{} \Distr-long=\Cn{} table and two \Pl-foot=\Cn{} \Distr-long=\Cn-3.\Ii, \Ccnj{} together five \Mass-3.\Ii \\ 
    \glt `The table is three feet long and two feet wide; together five [feet added together].'\hfill HH.\ElicitationDate{} (SF)
\end{exe}

HH gives a comment here: not \textit{sag̱ayt}, \textit{g̱asg̱oohl}. \textit{g̱asg̱oohl}, `equals'.

[1:07:30]
\begin{exe}
    \ex \glll
    gwile’lhl sese’ehl g̱a’negwihl  ii gilbilhl sese'’hl g̱ale'yhlit ha'niitx̱ooḵxw ii sag̱ayt k’ep g̱a-daaxhlit \\
    gwile’l=hl se-se’e=hl g̱a-’nakw=hl  ii gilbil=hl se-se’e=hl g̱a-le'y=hl-t ha’niitx̱ooḵxw ii sag̱ayt k’ep g̱a-daax=hl-t \\
    three=\Cn{} \Pl-foot=\Cn{} \Distr-long=\Cn{} \Ccnj{} three=\Cn{} \Pl-foot=\Cn{} \Distr-wide=\Cn-3.\Ii table \Ccnj{} together ten \Distr-perimeter=\Cn-3.\Ii \\
    \glt "." \hfill HH.\ElicitationDate{} (VF)
\end{exe}

HH: You could say \textit{g̱asg̱oo} if you really wanted to talk about the diameter, or the surrounding area. 

[1:10:37]
\begin{exe}
    \ex \glll
    k’ep g̱adaax g̱asg̱oohl ha’niitx̱ooḵxw \\
    k’ep g̱a-taax g̱asg̱oo=\Cn{} ha’niitx̱ooḵxw \\
    ten \Distr-perimeter \Mass-3.\Ii{} table \\
    \glt "The table's edge is ten in size; the table is ten around"  (orig. tr.) \hfill HH.\ElicitationDate{} (SF)
\end{exe}

[1:12:06]
\begin{exe}
    \ex \glll
    sag̱ayt ḵ’oolt g̱abiihl ts’ek 'nii t'aat lax̱ ha’niitx̱ooḵxw  \\
    sag̱ayt ḵ’oolt g̱abii=hl ts’ak 'nii t'aa-t lax̱ ha’niitx̱ooḵxw  \\
    together six \Cnt=\Cn{} plate on sit-3.\Sx{} on table \\
    \glt "In all six plates are on the table." \hfill HH.\ElicitationDate{} (VF)
\end{exe}

[1:14:24]
\begin{exe}
    \ex \glll
    gwile'lhl ts’ek ’nii t’aat lax̱ ha’niitx̱ooḵxw win t’aa ii gwile’lhl ts’ek t’aat andoosda \\
    gwile'l=hl ts’ek ’nii t’aa-t lax̱ ha’niitx̱ooḵxw win t’aa ii gwile’l=hl ts’ek t’aa-t an-toosta \\
    three=\Cn{} plate on sit-\Sx{} on table \Comp{} sit \Ccnj{} three=\Cn{} plate sit-\Sx{} \Nmlz-far.side \\
    \glt "There are three plates on this side of the table and three plates on the other side." \hfill HH.\ElicitationDate{} (XF)
\end{exe}

[1:14:44]
\begin{exe}
    \ex \glll
    gwile'lhl ts’ek t’aat win t’aa’a ii gwile’lhl andoosda \\
    gwile'l=hl ts’ek t’aa-t win t’aa-a'a? ii gwile’lhl andoosda \\
    three=\Cn{} plate sit-\Sx{} \Comp{} sit-\Detr? \Ccnj{} three=\Cn{} plate \Nmlz-far.side \\
    \glt "There are three plates [placed] and three on the other side." \hfill HH.\ElicitationDate{} (XF)
\end{exe}

It sounded like there was some glottalization toward the end of \textit{tcaa} between \textit{win} and \textit{ii}, not sure if that means there was something like -\textit{a’a} in there or not. 

[1:15:10]
\begin{exe}
    \ex \glll
    ’Wii ’nakwhl aats’ip tun \\
    ’wii ’nakw=hl aats’ip t=xʷin \\
    big long=\Cn{} door \Pn=\Dem.\Prox \\
    \glt "The door is tall." \hfill HH.\ElicitationDate{} (SF)
\end{exe}

HH corrects SF \textit{’nagwihl} to \textit{’nakwhl}. Cf. {g̱a’nagwihl}, where the vowel does seem to be added. 

[1:16:36]
\begin{exe}
    \ex \glll
    ’Wii ’nakwt Andrew (ii hlep dulpxw'y). \\
    ’wii ’nakw-t Andrew ii hlep dulp-xw-'y. \\
    big long=\Pn{} Andrew \Ccnj{} deep short-\Val-\Fsg.\Ii{} \\
    \glt "Andrew is tall (and I'm short)." \hfill HH.\ElicitationDate{} (XF)
\end{exe}

[1:17:14]
\begin{exe}
    \ex \glll
    ’Wii ’nakwt Andrew (ii hlep dulpxwin). \\
    ’wii ’nakw-t Andrew ii hlep dulp-xwin. \\
    big long=\Pn{} Andrew \Ccnj{} deep short-\Val-\Ssg.\Ii{} \\
    \glt "Andrew is tall (and you're short)." \hfill HH.\ElicitationDate{} (XF)
\end{exe}

[1:17:23]
\begin{exe}
    \ex \glll
    ’nakw win hlep dulpxwin es Andrew \\
    ’nakw win hlep dulp-xw-in a=s Andrew \\
    long \Comp{} deep short-\Val-\Ssg.\Ii{} \Prep=\Pn{} Andrew \\
    \glt "You are way shorter than Andrew." \hfill HH.\ElicitationDate{} (VF)
\end{exe}

Could try comparing this to \textit{dulp win hlep dulpxwin es Andrew} if someone is only a little shorter, etc.

[1:19:05]
\begin{exe}
    \ex \glll
    sag̱ayt k’i’yhl g̱asg̱oon \\
    sag̱ayt k’i’y=hl g̱asg̱oo-n \\
    together one=\Cn{} \Mass-2.\Ii \\
    \glt "We are the same size/height." \hfill HH.\ElicitationDate{} (VF)
\end{exe}

[1:19:10]
\begin{exe}
    \ex \glll
    sag̱ayt k’i’yhl g̱a’negwin \\
    sag̱ayt k’i’y=hl g̱a-’nakw-n \\
    together one=\Cn{} \Distr-long-2.\Ii \\
    \glt "We are the same height." \hfill HH.\ElicitationDate{} (VF)
\end{exe}

HH seems to be using -\textit{n} \Ssg.\Ii{} in plural contexts here (rather than -\textit{si'm} \Spl.\Ii).

[1:19:35]
\begin{exe}
    \ex \glll
    mooji sag̱ayt k’i’yhl g̱asg̱oo gant Andrew \\
    mooji sag̱ayt k’i’y=hl g̱asg̱oo gan=t Andrew \\
    almost together one=\Cn{} \Mass=\Pn{} Andrew \\
    \glt "[I am] almost the same height as Andrew." \hfill HH.\ElicitationDate{} (VF)
\end{exe}

Not sure where first person is in this, if at all; possible there could be a -\textit{'y} on \textit{g̱asg̱oo} that isn't quite audible?

[1:20:01]
\begin{exe}
    \ex[*]{\glll
    sag̱ayt k’i’yhl g̱a’negwi'y g̱an g̱a’negwit Andrew \\
    sag̱ayt k’i’y=hl g̱a-’nakw-'y g̱an g̱a-’nakw=t Andrew \\
    together one=\Cn{} \Distr-long-\Fsg.\Ii{} \Pcnj{} \Distr-long=\Pn{} Andrew \\
    \glt \textit{Intended: } "My height and Andrew's height are the same." \hfill HH.\ElicitationDate{} (SF)}
\end{exe}

[1:21:17]
\begin{exe}
    \ex[]{\glll
    sag̱ayt k’i’yhl g̱a’negwi'm g̱ant Andrew \\
    sag̱ayt k’i’y=hl g̱a-’nakw-'m g̱an=t Andrew \\
    together one=\Cn{} \Distr-long-\Fpl.\Ii{} \Pcnj=\Pn{} Andrew
    \\
    \glt "My height and Andrew's are the same; our heights are the same, [me] with Andrew." \hfill HH.\ElicitationDate{} (VF)}
\end{exe}

\begin{exe}
    \ex[]{\glll
    sag̱ayt k’i’yhl g̱a’negwi'm \\
    sag̱ayt k’i’y=hl g̱a-vnakw-'m \\
    together one=\Cn{} \Distr-long-\Fpl.\Ii{}
    \\
    \glt "Our heights are the same." \hfill HH.\ElicitationDate{} (VF)}
\end{exe}

[1:21:40]
\begin{exe} 
    \ex[]{\glll
    sag̱ayt k’i’yhl g̱a’negwi'y g̱ant Andrew \\
    sag̱ayt k’i’y=hl g̱a-’nakw-'y g̱an=t Andrew \\
    together one=\Cn{} \Distr-long-\Fsg.\Ii{} \Pcnj=\Pn{} Andrew
    \\
    \glt "My height and Andrew's are the same." \hfill HH.\ElicitationDate{} (VF)}
\end{exe}

[1:22:30]
\begin{exe}
    \ex[*]{\glll
    sag̱ayt k’i’yhl g̱a’nekwhl ha’niitx̱ooḵxw g̱anhl g̱ale’yhlit. \\
    sag̱ayt k’i’y=hl g̱a-’nakw=hl ha’niitx̱ooḵxw g̱anhl g̱a-le’y=hl-t. \\
    together one=\Cn{} \Distr-long=\Cn{} table \Pcnj=\Cn{} \Distr-long=\Cn-3.\Ii \\
    \glt \textit{Intended}: "The table is as long as it is wide." \hfill HH.\ElicitationDate{} (SF)}
\end{exe}

[1:23:07]
\begin{exe}
    \ex \glll
    sag̱ayt k’i’yhl g̱a’nekwhl ha’niitx̱ooḵxw g̱anhl g̱ale’ydiit. \\
    sag̱ayt k’i’y=hl g̱a-’nakw=hl ha’niitx̱ooḵxw g̱anhl g̱a-le’y-diit. \\
    together one=\Cn{} \Distr-long=\Cn{} table \Pcnj=\Cn{} \Distr-long-\Tpl.\Ii \\
    \glt "The table is as long as it is wide." \hfill HH.\ElicitationDate{} (VF)
\end{exe}

[1:24:37]
\begin{exe}
    \ex[*]{\glll
    ’wii ’nakwt Andrew as  'nii'y \\
    ’wii ’nakw=t Andrew a=s  'nii'y \\
    \Cntr{} more big long=\Pn{} Andrew \Prep=\Pn{} \Fsg.\Iii \\
    \glt \textit{Intended}: "Andrew is taller than I am." \hfill HH.\ElicitationDate{} (VF)}
\end{exe}

[1:24:45]
\begin{exe}
    \ex \glll
    g̱ay k’aa ’wii ’nakwt Andrew loo'y \\
    g̱ay k’aa ’wii ’nakw=t Andrew loo-'y \\
    \Cntr{} more big long=\Pn{} Andrew \Obl-\Fsg.\Ii \\
    \glt "Andrew is taller than I am." \hfill HH.\ElicitationDate{} (VF)
\end{exe}

The \textit{'y} of \textit{loo'y} was quite hard for me to hear here, but would be expected to appear.

\begin{exe}
    \ex \glll
    g̱ay k’aa ’wii ’nakwt Andrew loo'y, ii ḵ'ap nee dii  ’wii ’nakwt \\
    g̱ay k’aa ’wii ’nakw=t Andrew loo-'y, ii ḵ'ap nee dii  ’wii ’nakwt \\
    \Cntr{} more big long=\Pn{} Andrew \Obl-\Fsg.\Ii, \Ccnj{} \Verum{} \Neg{} \Foc{} big long-3.\Ii \\
    \glt "Andrew is taller than I am, but he's not really tall." \hfill HH.\ElicitationDate{} (SF)
\end{exe}

[1:25:40]
\begin{exe}
    \ex \glll
    g̱ay k’aa ’wii ’nakwt Andrew loo'y, ii ḵ'ap nee dii  ’wii ’nakw'm \\
    g̱ay k’aa ’wii ’nakw=t Andrew loo-'y, ii ḵ'ap nee dii  ’wii ’nakw-'m \\
    \Cntr{} more big long=\Pn{} Andrew \Obl-\Fsg.\Ii, \Ccnj{} \Verum{} \Neg{} \Foc{} big long-\Fpl.\Ii \\
    \glt "Andrew is taller than I am, but we're not very tall." \hfill HH.\ElicitationDate{} (VF)
\end{exe}

[1:27:53]
\begin{exe}
    \ex \glll
    k’i’yhl t’img̱es win ’wii ’nakw t’is \\
    k’i’y=hl t’img̱es win ’wii ’nakw t’is \\
    one=\Cn{} head \Comp{} big long big \\
    \glt "One head taller; taller by one head." \hfill HH.\ElicitationDate{} (VF)
\end{exe}

\textit{’wii ’nakw t’is} could alternatively be \textit{’wii ’nakwdis} with -\textit{t} 3.\Ii{} and =\textit{is} \Qudd. 